\lesson{1}{Nov 01 2021 Mon (08:16:07)}{The Legislative Branch}{Unit 3}

\subsubsection*{Bicameral Legislature}

\begin{definition}[Bicameral]
    Congress is \bf{\it{Bicameral}}. That means that it has two Houses. They are:
    
    \begin{itemize}
        \item The \bf{House of Representatives}
        \item The \bf{Senate}
    \end{itemize}
    
    For a bill to pass, it must pass the majority of both of the Houses of Congress to become a law. The \bf{House} gives more representation to states with larger populations. In contrast, state representation is equal in the \bf{Senate}, so smaller states have greater relative power there.
\end{definition}

As the nation grew, the number of \bf{House} members grew steadily until $1911$. That's when Congress set a fixed number of $435$ representatives. The \bf{U.S. Census Bureau} conducts an official count of a population carried out at a specific time every $10$ years. That's for the main purpose of apportionment. Each state must also have a representative.

The \bf{House} members are more expected than \bf{Senators} to be extra sensitive to the \it{State} and \it{Local} issues. The \bf{House} members are more dependent on support from constituents since they are \it{""Closer to the people"}. Each representative speaks for about $700,000$ people. But, senators reflect stability through their longer terms.

\subsubsection*{House}

\begin{itemize}
    \item At least one representative per state, based on population.
    \item Each member of the \bf{House} represents over a half-million people.
    \item $435$ members total.
    \item Two-year terms.
\end{itemize}

\subsubsection*{Senate}

\begin{itemize}
    \item Two senators per state, equal representation for states.
    \item $100$ members total.
    \item Six-year terms, one-third of the \bf{Senate} is up for re-election every two years.
\end{itemize}

\subsubsection*{Qualifications For Congress}

The \bf{House} and \bf{Senate} both have an age, citizenship, and residency requirements. Members of both the \bf{House} and \bf{Senate} must be residents of the state people elect them to represent. While the Constitution does not require the \bf{House} members to live in the specific district they represent, they usually do so by tradition.

Notice that the \bf{Senate} qualifications are higher for age and years of citizenship. The framers of the Constitution intended the \bf{Senate} to be the more stable \bf{House} of Congress with members likely to have greater political experience.

\subsubsection*{House}

\begin{figure}[h]
    \begin{quote}
        \it{"No Person shall be a Representative who shall not have attained to the Age of twenty five Years, and been seven Years a Citizen of the United States, and who shall not, when elected, be an Inhabitant of that State in which he shall be chosen."}
    \end{quote}
    \caption{From Article I, Section 2}
\end{figure}

These are the qualifications:

\begin{itemize}
    \item At least 25 years old
    \item A U.S. citizen for at least seven years
    \item Resident of the state chosen to represent
\end{itemize}

\subsubsection*{Senate}

\begin{figure}[h]
    \begin{quote}
        \it{"No Person shall be a Senator who shall not have attained to the Age of thirty Years, and been nine Years a Citizen of the United States, and who shall not, when elected, be an Inhabitant of that State for which he shall be chosen."}
    \end{quote}
    \caption{From Article I, Section 3}
\end{figure}

These are the qualifications:

\begin{itemize}
    \item At least 30 years old
    \item A U.S. citizen for at least nine years
    \item Resident of the state chosen to represent
\end{itemize}

\subsubsection*{Powers of Congress}

The Constitution specifies the areas where Congress may pass legislation in \bf{Article I, Section 8}. Some of these topics include collecting taxes and establishing the process for naturalization. Congress allocates funds collected to support specific programs and establishes agencies and rules to administer them.

The final statement in \bf{Article I, Section 8} is not a specific power. It states that:

\begin{figure}[h]
    \begin{quote}
        \it{"Congress has the power to make all laws which shall be necessary and proper for carrying into execution the foregoing powers, and all other powers vested by this Constitution in the government of the United States, or in any department or officer thereof."}
    \end{quote}
    \caption{From Article I, Section 8}
\end{figure}

We call this line the \it{"Elastic Clause"} because it \bf{"Stretches"} the powers of Congress to cover areas the framers of the Constitution could not have anticipated.

Interpreting the phrase \it{"Necessary and Proper"} has been controversial over the course of U.S. history. But, people debate what is \it{"Necessary and Proper"}. They also debate on what authority is implied by the powers listed in the Constitution. For example, since Congress has the power \it{"to raise and support an Army and Navy"}, officials decided the power implied the ability to create an air force. The framers could not have anticipated the invention of airplanes, but Congress assumed that the power to maintain a military force would certainly include the use of new technology and needs.

\subsubsection*{House}

\begin{itemize}
    \item All bills the will create revenue, meaning raise money through taxes, must originate in the \bf{House}.
    \item Only the \bf{House} can impeach high officials, like the president or a Supreme Court justice.
\end{itemize}

\subsubsection*{Senate}

\begin{itemize}
    \item The Senate must approve all foreign treaties made by the president.
    \item Only the Senate can conduct trials for and convict impeached federal officials.
    \item The Senate must confirm important appointments made by the president.
\end{itemize}

\subsubsection*{House VS Senate}

The structure of the \bf{House} and \bf{Senate} are different. The Constitution grants many specific powers which are shared between both of the houses, while it may just reserve some other powers for either of the houses.

The leadership in each house differs as well. The \bf{Speaker of the House} is the presiding officer of the \bf{House of Representatives}. Members of the political party holding a majority of the \bf{House} seats choose who will serve as the \bf{Speaker}. The \bf{Speaker} sets the agenda for the \bf{House}, determines who should serve on what committees, and what committees will study introduced bills. The \bf{Speaker} is second in the line of presidential succession. That means the Speaker would become acting president should the president and vice president become unable to perform their official duties.

The \bf{Senate's} presiding officer is technically the \bf{Vice President} according to the Constitution. However, this role has very little power, since the \bf{Vice President} can only vote on bills in the case of a tie in the \bf{Senate}. The \bf{President Pro Tempore}

However, this role has very little power. The reason is because the \bf{Vice President} can only vote on bills in the case of a tie in the \bf{Senate}. The \bf{President Pro Tempore} is the highest-ranking member of the \bf{Senate}. That happens only when the \bf{Vice President} is absent. And that happens a lot. Usually, the \bf{President Pro Tempore} is the member of the majority political party who has served in the \bf{Senate} the longest. The role of the \bf{President Pro Tempore} is similar to the \bf{Speaker of the House}. The \bf{President Pro Tempore} is third in the line of succession after the Speaker of the House. Both the \bf{Senate} and the \bf{House} give a leadership position to a member of the political party in the minority as well.

\subsubsection*{Congress Making Laws}

\begin{figure}[h]
    \begin{quote}
        \it{"Each House may determine the Rules of its Proceedings… (and) shall keep a Journal of its Proceedings, and from time to time publish the same."}
    \end{quote}
    \caption{From Article I, Section 5}
\end{figure}

Each house keeps a record of its actions and debates. You can access current and historical actions of Congress on the Internet.

Congress determines many specifics of the lawmaking process for itself. But, it takes guidance from the Constitution. It can determine when to hold a vote on a bill and when to send it to another group of Congress members for further study and edits. A majority of both the \bf{House} and \bf{Senate} must pass new legislation and send it to the \bf{President} for approval to become a law.

The \bf{House} conducts its business more quickly than the \bf{Senate} and has limits on length of debate for a bill. The \bf{Senate} has no limit on debate. And even sometimes, a senator will intentionally delay voting on a bill by talking as long as possible or reading a long piece of text. This is called a \bf{Filibuster}. A vote of cloture can end a \bf{Filibuster} and force a vote, with at least $60$ senators in agreement.

\subsubsection*{Law Making Process}

\paragraph{Bill Introduction}

\bf{House}: A bill is introduced in the \bf{House}. It's given a number and referred to a committee.
\bf{Senate}: A bill is introduced in the \bf{Senate}. It's given a number and referred to a committee.

All of the bills start as an idea. Bills may begin in either the \bf{House of Representatives} or the \bf{Senate}. The representative or senator introduces his or her bill. The idea is read, given a number, and referred to a committee or subcommittee. The \bf{"H", "S"} in front of a bill number tells you whether it began in the \bf{House} or the \bf{Senate}. About $9,000$ bills are introduced each year in Congress, but only about $10\%$ become laws.

\paragraph{Subcommittee Review}

\bf{House}: A \bf{House} subcommittee often first considers the bill. It votes on advancing the bill to the full committee.
\bf{Senate}: A \bf{Senate} subcommittee often first considers the bill. If votes on advancing the bill to the full committee.

Most of the work on bills happens in committees. The legislators might edit the language, make additions, or take things out. Together, the \bf{House} and \bf{Senate} have $40$ committees and $160$ subcommittees. Representatives and senators will serve on committees or subcommittees that match their background or areas of interest.

\paragraph{Committee Review}

\bf{House}: A full committee in the \bf{House} considers and may change the bill. Then it will vote on advancing the bill.
\bf{Senate}: A full committee in the \bf{Senate} considers and may change the bill. Then it vote on advancing the bill.

Work in the full committee is the hardest step for a bill to pass. Only $1$ in $6$ bills will pass this step. The full committee may make more edits to the bill's language before taking a vote. If the bill does pass, it will go to the whole \bf{House} or \bf{Senate} for debate. For the \bf{House}, first they setup rules about the debate in the \bf{House Rules Committee}.

\paragraph{House Rules Committee}

\bf{House}: The \bf{House Rules Committee} may make special rules for debating and changing the bill. The \bf{House} votes on the rules.

The \bf{House Rules Committee} will set rules for debate on the bill. Examples of these rules can include the time allowed for debate. It might limit the number of changes that can be made to the bill. The \bf{Senate} has fewer rules and unlimited time for debate.

\paragraph{Floor Action}

\bf{House}: The \bf{House} debates and may change the bill. Members vote on its passage.
\bf{Senate}: The \bf{Senate} debates and may change the bill. Members vote on its passage.

Debate of the whole \bf{House} or \bf{Senate} is called the \it{"Floor Action"}. This step is challenging because representatives and senators come from different parts of the nation. Each member has different interests and concerns. The bill needs approval from a majority of people in both chambers to advance.

\paragraph{House-Senate Compromise}

\bf{House}: \bf{House} and \bf{Senate} compromise on differences between each version of the bill.
\bf{Senate}: \bf{House} and \bf{Senate} compromise on differences between each version of the bill.

A conference committee with members from both the \bf{House} and \bf{Senate} meet. They discuss differences between each chamber's version of the bill. They will compromise on the differences to create a single bill. A compromise happens when both sides give up some things they originally wanted.

\paragraph{Final House, Senate Vote}

\bf{House}: The \bf{House} debates the compromise bill and votes on it.
\bf{Senate}: The \bf{Senate} debates the compromise bill and votes on it.

The compromise bill is sent back to both the \bf{House} and the \bf{Senate}. Members in each debate the compromise version of the bill. They vote on whether to approve it. Approval from both the \bf{House} and the \bf{Senate} is required on the same exact version of the bill for it to continue.

\subsubsection*{Committees Organization}

Committees do most of the work of lawmaking and have great power in determining whether a bill continues or dies. In fact, most bills originate from the work of committees. Members of committees research issues and potential solutions, debate ideas, and discuss the interests of different groups of people affected by the issues and bills.

Usually people who have seniority in Congress serve on the major committees. Generally, members of a committee have a background or special interest in the topic focus of the committee.

Standing–Bills are usually referred to a standing committee. They are permanent and have a specific area of responsibility.

\bf{Ad hoc–An ad hoc} committee is temporary for the purpose of creating or enforcing mandates for a specific reason. Congress disbands an ad hoc committee when the need is fulfilled.

\bf{Joint–A joint} committee consisting of members from both the \bf{House} and the \bf{Senate} forms to address important issues needing special attention from both. \bf{Joint committees} make unified recommendations to the whole Congress.

\bf{Conference–A conference} committee has members from both houses and forms to settle differences between similar bills passed separately in the House and Senate.

\bf{House Rules Committee–This committee}, specific to the \bf{House}, sets rules of debate for particular bills.

\bf{The Joint Economic Committee (JEC)} studies and makes recommendations to the whole Congress regarding national economic conditions. A committee hearing is not a criminal or civil trial proceeding like in the judiciary system. Instead, congressional hearings are public meetings intended to gather information to inform decision on policy and law.

\subsubsection*{Influences on Congress Member's Vote}

When a member of Congress takes office, they pledge to represent the interests of the people from their district or state. However, even within a small district, ideas and concerns can be very diverse. Remember, each House member represents over a half-million people! So how do they decide how to vote on bills?

Many factors influence the decisions made by members of Congress. Lawmakers study and consider their constituents' desires concerning proposed legislation. If polls and communications indicate that a majority of constituents want a bill to pass, lawmakers may be more likely to vote in favor of it. Lawmakers' own values and beliefs also influence how they vote on bills. Often they must depend on their own opinions, especially when their constituents are deeply divided on an issue. Interest groups and political parties also seek to sway members of Congress to vote for or against particular bills. Ultimately, the individuals who represent us in Congress have the difficult job of considering all the information and opinions available and making the best decision for their constituents. The amount of influence constituents, personal opinions, or organizations like interest groups and political parties have on a lawmaker's vote varies by issue and differs for each member of Congress.

\bf{Republican Senator Susan Collins} from Maine was interviewed by reporters before an initial vote on a new economic investment bill. The bill met resistance, as is typical in congressional negotiations. However, a committee made up of both Democrats and Republicans, including Collins, put together a bill that passed in the Senate. The House presented its own stipulations for the bill, which is also normal in policy making. The multiple steps and delays in considering new bills helps make sure the process and final bill is thorough. Representatives must serve their constituents and consider the long-term impact of their decisions. Though sometimes the public thinks Congress moves too slowly, debate, compromise, and time are essential ingredients in creating new laws.

\newpage
